\section{Conclusion}

\paragraph{Summarisation of the project.}

As per the introduction (section~\ref{sec:introduction}), the aim of this
project was to implement a suspicious luggage detection system in an embedded
environment to bridge the gap between academic research and a real-world edge
application. The implementation proves the viability of devices such as the
NVIDIA Jetson Orin Nano for this task, demonstrating that a deployable solution
that can integrate with surveillance infrastructure is achievable within the
constraints of an embedded system.

Implementation of the system took the form of an end-to-end pipeline, utilising
a fine-tuned YOLO26s model for object detection, NvDeepSORT tracking with a
ResNet50-based re-identification model, custom spatio-temporal abandonment
logic deployed on a Jetson Orin Nano via NVIDIA's DeepStream SDK. Multiple
iterations found that a two-phase training strategy was optimal, with a COCO
\citep{lin2014microsoft} and OID \citep{krasin2017openimages} base, followed by
fine-tuning on the domain-specific ABODA dataset \citep{lin2015aboda}.

Evaluation of the model was completed on multiple datasets, utilising a custom
evaluation framework featuring event-based ground truth labelling.

\subsection{Results and Limitations}

The object detection model achieved a mAP@0.5 of $0.859$ on the ABODA test set,
balanced across the two classes of ``luggage'' and ``person''. The pipeline
performed well on in-domain data, achieving an F1-score of $0.774$ on the ABODA
examples with a precision of $0.750$ and recall of $0.800$.

However, out-of-domain performance on the IITP20 dataset revealed a significant
generalisation gap, with F1-scores ranging from $0.265$ to $0.346$ across the
three difficulty tiers. Precision was consistently the weakest metric across
all datasets, indicating that the system over-triggers rather than
under-detects. Three recurring failure modes were identified: false ownership
assignment of pre-existing stationary luggage, tracker ID switching in crowded
or occluded scenes, and difficulty re-triggering after an initial abandonment
detection in multi-event sequences.

From a hardware perspective, the Jetson Orin Nano sustained real-time
processing across all evaluation runs without thermal throttling, with the GPU
operating at $86$--$90\%$ average utilisation. This confirms the viability of
the platform for the target workload, though the high GPU utilisation leaves
limited headroom for additional processing or more complex models.

\subsection{Reflection}

The gap between in-domain and out-of-domain performance highlights the
difficulty of generalising a surveillance pipeline trained on a limited set of
environments. The ABODA results, while strong, were shown to be inflated by the
training domain overlap --- a limitation that was identified and documented
during the evaluation process. Relying on the IITP20 and AVSS results as the
more honest measure of performance, the system in its current form would not be
suitable for deployment without further refinement to address the false
positive rate.

The decision to include ABODA in the fine-tuning data improved detection
accuracy in the target domain but came at the cost of evaluation independence.
In hindsight, reserving ABODA exclusively for evaluation and sourcing
additional CCTV-style training data from other datasets would have provided a
cleaner separation between training and evaluation, and a more confident
measure of the system's true performance.

The evaluation framework itself --- event-based ground truth with temporal
tolerance --- proved to be a practical approach for assessing abandonment
detection without requiring frame-level annotations. This methodology could be
applied to future work in the domain.

\subsection{Future Work}

Several areas for improvement were identified through the evaluation:

\begin{itemize}
	\item \textbf{Scene initialisation.} A warm-up period during which stationary
	      objects are catalogued and excluded from ownership assignment would
	      address the dominant source of false positives observed in the IITP20
	      dataset.
	\item \textbf{Tracker robustness.} The zero-IoU true positive pattern
	      observed across multiple IITP20 videos points to tracker ID switching as
	      a systematic issue. A more robust re-identification model, or a tracker
	      architecture better suited to occluded environments, would improve both
	      precision and spatial accuracy.
	\item \textbf{Custom tracking implementation.} The current reliance on NVIDIA's
	      ReID model and DeepSORT architecture, while convenient, may not be optimal for the
	      specific challenges of the domain. A custom tracking solution, designed with the
	      unique requirements of suspicious luggage detection in mind, could yield better
	      performance and efficiency.
	\item \textbf{Multi-event handling.} Investigating why the pipeline fails to
	      reset after an initial detection would improve recall on sequences
	      containing multiple abandonment events.
	\item \textbf{Training data diversity.} Expanding the training dataset with
	      more CCTV-native footage from varied environments --- particularly
	      crowded scenes, overhead angles, and cluttered backgrounds --- would
	      help close the generalisation gap observed between ABODA and IITP20.
	\item \textbf{Model optimisation.} Exploring INT8 quantisation would reduce
	      the GPU utilisation overhead, freeing headroom for additional processing
	      tasks or enabling deployment on lower-powered hardware.
\end{itemize}

\paragraph{Final thoughts.}

This project demonstrates the feasibility of an embedded suspicious luggage
detection system, while also highlighting the significant challenges that
remain in achieving robust, generalisable performance in real-world
surveillance environments. The results provide a strong foundation for future
work, with clear avenues for improvement identified through the evaluation
process. With further refinement and additional training data, a deployable
solution that effectively balances precision and recall could be within reach.
